\subsection{EUCLIDEAN DIVISION}

THEOREM 1. Let \(a\) and \(b\) be integers such that \(b > 0\) . Then there exist integers \(q\) and \(r\) such that \(a = bq + r\) and \(r < b\) , and the integers \(q\) and \(r\) are uniquely determined by these conditions.

The conditions on \(q\) and \(r\) are equivalent to \(bq \leqslant a < b(q + 1)\) and \(r = a - bq\) (no. 2, Proposition 2). Hence we have to find \(q\) such that \(bq \leqslant a < b(q + 1)\) ; in other words, \(q\) must be the smallest integer such that \(a < b(q + 1)\) , which shows that \(q\) and \(r = a - bq\) are uniquely determined. To prove their existence, we note that there exist integers \(p\) such that \(a < bp\) , for example \(a + 1\) (since \(b > 0\) ). Let \(m\) be the least of these integers. We have \(m \neq 0\) , and we may therefore write \(m = q + 1\) with \(q \leqslant m\) ( \(\S 4\) , no. 2, Proposition 2); it follows that \(bq \leqslant a < b(q + 1)\) .

DEFINITION 1. With the notation of Theorem 1, r is called the remainder of the division of a by b. If r = 0, we say that a is a multiple of b, or that a is divisible by b, or that b is a divisor of a, or that b divides a, or that b is a factor of a. The number q is then called the quotient of a by b and is denoted by \(\frac{a}{b}\) or a/b.

If \(a\) is not a multiple of \(b\) , the number \(q\) is called the integral part of the quotient of \(a\) by \(b\) (cf.~General Topology, Chapter IV, §8, no. 2).

In this chapter, writing \(a / b\) or \(\frac{a}{b}\) will imply that \(b\) divides \(a\) .

The relations \(a = bq\) and \(q = a / b\) are equivalent (if \(b > 0\) ). Every multiple \(a'\) of a multiple \(a\) of \(b\) is a multiple of \(b\) , and

\[
a ^ {\prime} / b = (a ^ {\prime} / a) (a / b) \quad \text { if } \quad a \neq 0.
\]

Also, if \(c\) and \(d\) are multiples of \(b\) , then \(c + d\) and \(c - d\) (if \(d \leqslant c\) ) are multiples of \(b\) , and we have

\[
\frac {c + d}{b} = \frac {c}{b} + \frac {d}{b}, \quad \frac {c - d}{b} = \frac {c}{b} - \frac {d}{b}.
\]

The integers which are multiples of 2 are said to be even, and the others odd. By Theorem 1 the odd integers are of the form \(2n + 1\) .

